% ==========================================================
\section{防控策略优化}
% =========================================================

\subsection{干预措施的类型与效果}

基于敏感性分析的结果，本文将七类干预措施参数化，量化其对传播动力学的影响．这些措施可分为三类：

\textbf{第一类：降低传播系数的措施}

\begin{itemize}
  \item \textbf{口罩佩戴}（$u_1$）：外科口罩对流感病毒的过滤效率约17\%—25\%，N95口罩可达95\%以上\textsuperscript{[10]}．考虑到校园场景中学生多使用外科口罩，本文取保守估计：口罩佩戴率为$u_1$时，传播系数降低为$\beta_0 \times (1 - 0.17 u_1)$．当$u_1 = 1$（全员强制口罩）时，$\beta_0$降低17\%．
  
  \item \textbf{教室通风}（$u_2$）：增加教室换气次数能够稀释空气中的病毒浓度，降低气溶胶传播风险．文献报道，将换气次数从3次/小时提升至6次/小时，可降低传播风险约22\%\textsuperscript{[10]}．本文将通风改善精细化作用于宿舍与教室的$\beta$修饰系数，按室内场所比例加权分配，当$u_2 = 1$时，室内传播系数等效降低约20\%．
  
  \item \textbf{环境消毒}（$u_7$）：对高频接触表面（门把手、桌面、电梯按钮等）进行定期消毒，可降低接触传播风险．然而，流感的主要传播途径为飞沫和气溶胶，接触传播的贡献较小，消毒的效果有限，取$\beta_0 \times (1 - 0.05 u_7)$，最大降低5\%．
\end{itemize}

\textbf{第二类：减少接触的措施}

\begin{itemize}
  \item \textbf{线上授课}（$u_5$）：将部分或全部课程转为线上，可消除教室场景的接触．根据场所权重设定，教室接触占总接触的25\%，故线上授课可降低总接触率约25\%．参数化：当$u_5 = 1$时，接触矩阵$C_{ij}$中教室场景的权重归零，总接触率降低25\%．
  
  \item \textbf{社团限流}（$u_6$）：暂停或限制社团活动、体育锻炼、大型聚会等课外聚集，可减少课外场景的接触．课外接触占总接触的45\%，若限制程度为$u_6$，则课外接触降低$0.7 u_6$（最大降低70\%，因为部分课外活动无法完全取消）．参数化：课外接触权重调整为$0.45 \times (1 - 0.7 u_6)$，当$u_6 = 1$时，课外接触降低约24\%（$0.45 \times 0.7 = 0.315$，降低$0.45 - 0.315 = 0.135$，相对降幅$0.135/0.45 = 30\%$，占总接触的$0.45 \times 0.7 = 0.315$，即31.5\%）．
\end{itemize}

\textbf{第三类：阻断传播链的措施}

\begin{itemize}
  \item \textbf{疫苗接种}（$u_3$）：提高疫苗接种覆盖率，可增加免疫人群比例，减少易感者数量．参数化：接种率调整为$v \leftarrow v + u_3 \cdot \Delta v$，其中$\Delta v$为接种率提升幅度，目标上限取50\%（即从0.32\%最高提升至50\%）．当$u_3 = 0.8$时，接种率提升至约40\%，能够接近群体免疫阈值（\textbf{12.5\%}）．
  
  \item \textbf{病例隔离}（$u_4$）：加强症状监测和隔离响应，缩短从症状出现到隔离的时间延迟，提高隔离率$\alpha$．参数化：隔离率调整为$\alpha \times (1 + 3 u_4)$，当$u_4 = 1$时，隔离率提升4倍，能够有效缩短症状出现至隔离的时间延迟．
\end{itemize}

七类干预措施的参数化汇总见表\ref{tab:interventions}：

\begin{table}[H]
  \centering
  \caption{七类干预措施参数化方案}
  \label{tab:interventions}
  \small
  \begin{tabular}{cclcc}
    \toprule
    编号 & 变量 & 措施 & 作用参数 & 最大效果 \\
    \midrule
    $u_1$ & 口罩佩戴 & 外科口罩强制 & $\beta_0 \times (1-0.17u_1)$ & 降低17\% \\
    $u_2$ & 教室通风 & 换气次数提升 & $\beta_{\text{dorm/class}} \times \left(1 - 0.40\, u_2 \cdot w_v\right)$ & 降低20\% \\
    $u_3$ & 疫苗接种 & 季节性接种活动 & $v \leftarrow v+u_3\Delta V$ & 覆盖率+50\% \\
    $u_4$ & 病例隔离 & 隔离响应强化 & $\alpha \times (1+3u_4)$ & 提升4倍 \\
    $u_5$ & 线上授课 & 全线上切换 & $C_{ij}$教室权重降低  & 降低25\% \\
    $u_6$ & 社团限流 & 课外接触削减 & $C_{ij}$户外$\times(1-0.35u_6)$ & 降低35\% \\
    $u_7$ & 环境消毒 & 每日消毒 & $\beta_0 \times (1-0.05u_7)$ & 降低5\% \\
    \bottomrule
  \end{tabular}
\end{table}

\subsection{优化目标函数的构建}

\subsubsection{多目标优化的必要性}

防控策略的制定需要在多个目标间权衡：

\begin{itemize}
  \item \textbf{最小化感染规模}：降低感染人数，保护师生健康；
  \item \textbf{最小化防控成本}：减少口罩、消毒、线上教学设备等经济投入；
  \item \textbf{最小化教学影响}：减少线上授课对教学质量的负面影响．
\end{itemize}

这三个目标往往是冲突的：全面线上授课（$u_5=1$）能最大程度降低感染，但对教学影响巨大；完全不干预（所有$u_k=0$）成本最低，但感染规模最大．因此，需要构建多目标优化模型，在三个目标间寻求平衡．

\subsubsection{目标函数的定义}

本文采用加权和法，将三个目标整合为单一目标函数：

\begin{equation}
  \min F = w_1 \cdot C_{\text{infect}} + w_2 \cdot C_{\text{cost}} + w_3 \cdot C_{\text{edu}}
  \label{eq:objective}
\end{equation}

各分量的定义：

\textbf{感染成本}$C_{\text{infect}}$：累计感染总人数，归一化至$[0, 1]$区间：

\begin{equation}
  C_{\text{infect}} = \frac{R_{\text{total}}(u)}{R_{\text{total}}^{\text{baseline}}}
\end{equation}

其中，$R_{\text{total}}(u)$为干预$u$下的累计感染总人数，$R_{\text{total}}^{\text{baseline}}$为基准情景的累计感染总人数．$C_{\text{infect}} = 1$表示与基准情景相同，$C_{\text{infect}} = 0$表示零感染（理想情况）．

\textbf{防控成本}$C_{\text{cost}}$：综合各类措施的经济投入，归一化至$[0, 1]$：

\begin{equation}
  C_{\text{cost}} = \sum_{k=1}^{7} c_k \cdot u_k
\end{equation}

其中，$c_k$为措施$k$的单位成本系数（假设最大成本为1，总成本$\leq 7$，归一化至$[0, 1]$）．各措施的成本系数设定：

\begin{itemize}
  \item 口罩佩戴（$c_1 = 0.10$）：口罩单价低，但需大量采购；
  \item 教室通风（$c_2 = 0.15$）：设备采购和运行成本；
  \item 疫苗接种（$c_3 = 0.20$）：疫苗采购和接种组织成本；
  \item 病例隔离（$c_4 = 0.10$）：隔离场所和管理成本；
  \item 线上授课（$c_5 = 0.25$）：平台、设备、培训成本；
  \item 社团限流（$c_6 = 0.10$）：管理成本；
  \item 环境消毒（$c_7 = 0.10$）：消毒用品和人力成本．
\end{itemize}

\textbf{教学影响}$C_{\text{edu}}$：主要反映线上授课的影响，归一化至$[0, 1]$：

\begin{equation}
  C_{\text{edu}} = u_5
\end{equation}

即教学影响仅与线上授课的比例$u_5$相关，$u_5 = 0$时无影响，$u_5 = 1$时影响最大．这一简化忽略了其他措施对教学的影响（如社团限流影响课外活动），但抓住了主要矛盾．

\textbf{权重系数}$w_1, w_2, w_3$：反映决策者对各目标的重视程度，满足$w_1 + w_2 + w_3 = 1$．本文考虑三套权重方案：

\begin{itemize}
  \item \textbf{最小感染优先}：$(w_1, w_2, w_3) = (0.6, 0.2, 0.2)$，适用于聚集性疫情场景，首要目标是控制感染；
  \item \textbf{最小成本优先}：$(w_1, w_2, w_3) = (0.2, 0.6, 0.2)$，适用于常态化防控，资源约束较强；
  \item \textbf{最小教学影响优先}：$(w_1, w_2, w_3) = (0.2, 0.2, 0.6)$，适用于教学质量优先的场景．
\end{itemize}

\subsection{优化问题的形式化}

综合以上定义，优化问题形式化为：

\begin{align}
  \min_{u_1, \ldots, u_7} \quad & F = w_1 \cdot \frac{R_{\text{total}}(u)}{R_{\text{total}}^{\text{baseline}}} + w_2 \cdot \sum_{k=1}^{7} c_k u_k + w_3 \cdot u_5 \\
  \text{s.t.} \quad & 0 \leq u_k \leq 1, \quad k = 1, \ldots, 7
\end{align}

这是一个7维的约束优化问题，目标函数$F$通过求解ODE模型获得，无解析表达式，需采用数值优化方法求解．

\subsection{求解方法}

由于目标函数无解析表达式且可能存在多个局部最优，本文采用\textbf{网格搜索}（Grid Search）方法，在每个控制变量的取值范围内离散采样，遍历所有组合，找到最优解．

具体步骤：

\begin{enumerate}
  \item 对每个控制变量$u_k$在$[0, 1]$区间内离散为11个取值：$\{0, 0.1, 0.2, \ldots, 1.0\}$；
  \item 遍历所有$11^7 \approx 1.9 \times 10^7$种组合；
  \item 对每种组合，求解ODE模型，计算目标函数$F$；
  \item 找到$F$最小的组合作为最优解．
\end{enumerate}

\subsection{分场景最优方案}

本文定义三种场景，分别优化防控策略：

\textbf{常态场景}：无聚集性疫情信号，日常防控状态．优化结果：

\begin{table}[H]
  \centering
  \caption{常态场景最优防控方案}
  \begin{tabular}{c|ccccccc|c}
    \toprule
    季节 & $u_1$ & $u_2$ & $u_3$ & $u_4$ & $u_5$ & $u_6$ & $u_7$ & 效果 \\
    \midrule
    冬季 & 0.3 & 0.5 & 0.8 & 0.2 & 0 & 0.3 & 0.5 & 降低25\% \\
    夏季 & 0.2 & 0.5 & 0.8 & 0.2 & 0 & 0.2 & 0.5 & 降低37\% \\
    \bottomrule
  \end{tabular}
\end{table}

\textbf{解读}：常态防控以疫苗接种推广（$u_3=0.8$）和教室通风（$u_2=0.5$）为主，无需全员口罩和线上授课．疫苗接种率从0.32\%提升至约16\%，虽然仍低于群体免疫阈值，但能提供基础保护．冬季的防控强度略高于夏季，反映了冬季峰的高风险．

\textbf{散发场景}：出现零星确诊病例，需加强防控．优化结果：

\begin{table}[H]
  \centering
  \caption{散发场景最优防控方案}
  \begin{tabular}{c|ccccccc|c}
    \toprule
    季节 & $u_1$ & $u_2$ & $u_3$ & $u_4$ & $u_5$ & $u_6$ & $u_7$ & 效果 \\
    \midrule
    冬季 & 1.0 & 1.0 & 0.8 & 1.0 & 0 & 1.0 & 1.0 & AR降低\textbf{64\%} \\
    夏季 & 1.0 & 1.0 & 0.8 & 1.0 & 0 & 0.5 & 1.0 & AR降低\textbf{41\%} \\
    \bottomrule
  \end{tabular}
\end{table}

\textbf{解读}：散发场景须全面启动口罩强制（$u_1=1.0$）、教室通风（$u_2=1.0$）、病例强隔离（$u_4=1.0$）及社团停摆（$u_6=1.0$），线上授课不启动（$u_5=0$，以避免接触向宿舍集中）．预期累计感染降低约\textbf{41\%—64\%}（夏季至冬季），隔离率提升4倍，有效阻断传播链．

\textbf{聚集场景}：出现班级或宿舍连片感染，需全面防控．优化结果：

\begin{table}[H]
  \centering
  \caption{聚集场景最优防控方案}
  \begin{tabular}{c|ccccccc|c}
    \toprule
    季节 & $u_1$ & $u_2$ & $u_3$ & $u_4$ & $u_5$ & $u_6$ & $u_7$ & 效果 \\
    \midrule
    冬季 & 1.0 & 0.8 & 0.8 & 1.0 & 1.0 & 1.0 & 1.0 & 降低43\%—69\% \\
    夏季 & 1.0 & 0.8 & 0.8 & 1.0 & 0.8 & 1.0 & 1.0 & 降低36\%—63\% \\
    \bottomrule
  \end{tabular}
\end{table}

\textbf{解读}：聚集场景全面启动所有措施：全员口罩（$u_1=1.0$）、全线上授课（$u_5=1.0$）、社团停摆（$u_6=1.0$）、强隔离（$u_4=1.0$），峰值感染可降低约\textbf{43\%}，累计感染可降低约\textbf{63\%—69\%}（夏季至冬季）．这一方案虽然成本和教学影响较大，但能够有效遏制聚集性疫情的蔓延，避免更大规模的爆发．

\subsection{优化结果的验证}

为验证优化结果的合理性，本文对比了不同干预组合的效果，确认最优解的有效性．图\ref{fig:compare}展示了部分对比结果：

\begin{figure}[H]
  \centering
  \includegraphics[width=0.9\textwidth]{photos/06_scenario_comparison.png}
  \caption{不同场景下防控方案效果对比}
  \label{fig:compare}
\end{figure}

对比结果表明：聚集场景的全面防控方案能够将累计感染人数从基准的\textbf{28215}人降低至\textbf{8838}人，降幅约\textbf{69\%}（峰值感染率降低约43\%）；散发场景的全面防控能够使累计感染降低约\textbf{41\%—64\%}；常态场景的轻度防控能够将峰值感染率降低约\textbf{25\%—37\%}，符合预期．
